\section{What a semidefinite representation would imply}
\label{sec:pencil}

The two slices in Proposition~\ref{prop:slices} constrain every pencil
representing $K$.  We now turn those constraints into an exact block
matrix inequality whose off-diagonal block is linear in $y$.
A positive linear map preserves positive semidefiniteness; it is
called \emph{unital} if it sends the identity to the identity.

\begin{proposition}\label{prop:pencil}
Let $K=\Lambda_+(p,(I_4,I_4,0))$ for the polynomial
in~\eqref{eq:polynomial}. Suppose that $K$ admits a representation by a finite homogeneous real
symmetric linear pencil.  Then there are integers $a,c\geq1$, positive
linear maps
\[
 D:\Sym^4\longrightarrow\Sym^a,
 \qquad E:\Sym^4\longrightarrow\Sym^c,
 \qquad D(I_4)=I_a,\quad E(I_4)=I_c,
\]
and a linear map $B:\R^3\longrightarrow\R^{a\times c}$ such that, for
every $X\succ0$, $Z\in\Sym^4$ and $y\in\R^3$,
\begin{equation}\label{eq:schur-equivalence}
 Z\succeq\Phi_y(X^{-1})
 \quad\Longleftrightarrow\quad
 E(Z)\succeq B(y)^\top D(X)^{-1}B(y).
\end{equation}
\end{proposition}

\begin{proof}
Write the hypothetical pencil as
\[
 L(X,Z,y)=L_1(X)+L_2(Z)+L_3(y).
\]
We first remove the pencil's common kernel and normalize
two positive blocks.
We then rescale the pencil and prove that its limit represents the
same inequality.

\emph{Compression and normalization.}
We use the elementary fact that
\[
 A+H\succeq0,\quad A-H\succeq0
 \quad\Longrightarrow\quad \ker A\subseteq\ker H.
\]
Indeed, if $w\in\ker A$, the two nonnegative quadratic forms on $w$
sum to zero.  A positive semidefinite matrix annihilates every vector
on which its quadratic form vanishes, so both $(A+H)w$ and $(A-H)w$
are zero.

Since $e$ is interior to $K$, for every ambient vector $x$ there is
an $\varepsilon>0$ with $L(e)\pm\varepsilon L(x)\succeq0$.
The preceding fact shows that $\ker L(e)$ is a common kernel of the
whole pencil.  Compressing to its orthogonal complement preserves the
represented cone and makes $L(e)\succ0$.  This complement is nonzero:
otherwise the pencil would vanish identically, whereas
$(-I_4,I_4,0)\notin K$ by~\eqref{eq:slice-zero}.

The same slice shows that $L_1$ and $L_2$ are positive maps.  Set
$U=\ker L_1(I_4)$.  Positivity at $I_4\pm\varepsilon X$, for small
$\varepsilon>0$, and the same kernel argument show that every
$L_1(X)$ annihilates $U$.  Thus, in $U^\perp\oplus U$,
\[
 L_1(X)=\begin{pmatrix}D(X)&0\\0&0\end{pmatrix}.
\]
Both summands are nonzero.  If $U^\perp=0$, then $L_1=0$ and the
pencil accepts $(-I_4,I_4,0)$, contradicting~\eqref{eq:slice-zero}.
If $U=0$, then $L_1(I_4)\succ0$, so
$L_1(I_4)-\varepsilon L_2(I_4)\succ0$ for sufficiently small
$\varepsilon>0$.  This would accept $(I_4,-\varepsilon I_4,0)$,
again contradicting that slice.

Let $a=\dim U^\perp$, $c=\dim U$, and let $E(Z)$ be the lower
diagonal block of $L_2(Z)$.  Both $D$ and $E$ are positive maps.
Moreover, $D(I_4)\succ0$ by the definition of $U$, and
$E(I_4)\succ0$ because it is the compression of $L(e)\succ0$ to
$U$.  Apply the fixed block diagonal congruence
\[
 \operatorname{diag}\bigl(D(I_4)^{-1/2},E(I_4)^{-1/2}\bigr)
\]
and retain the notation for the resulting blocks.  We now have
$D(I_4)=I_a$ and $E(I_4)=I_c$.

\emph{Rescaling the pencil.}
Write the other blocks as
\[
 L_2(Z)=\begin{pmatrix}F(Z)&C(Z)\\C(Z)^\top&E(Z)\end{pmatrix},
 \qquad
 L_3(y)=\begin{pmatrix}A(y)&B(y)\\B(y)^\top&G(y)\end{pmatrix}.
\]
Fix $y$ and put $Z_0=\Phi_y(I_4)$.  By
\eqref{eq:slice-positive} and $\Phi_{sy}=s^2\Phi_y$, the point
$(I_4,s^2Z_0,sy)$ lies in $K$ for every real $s$.  Its lower pencil
block therefore satisfies
\[
 s^2E(Z_0)+sG(y)\succeq0 \qquad(s\in\R).
\]
Dividing by $|s|$ and taking limits through positive and negative
$s$ gives $G(y)\succeq0$ and $-G(y)\succeq0$.  Hence $G(y)=0$.

For fixed $X\succ0$, $Z$ and $y$, and any real $s\ne0$, congruence
of $L(X,s^2Z,sy)$ by $\operatorname{diag}(I_a,s^{-1}I_c)$ gives
\begin{equation}\label{eq:scaled-pencil}
 M_s(X,Z,y)=
 \begin{pmatrix}
 D(X)+sA(y)+s^2F(Z)&B(y)+sC(Z)\\
 B(y)^\top+sC(Z)^\top&E(Z)
 \end{pmatrix}.
\end{equation}
Congruence preserves positive semidefiniteness for either sign of
$s$.  The exact slice~\eqref{eq:slice-positive} gives
\begin{equation}\label{eq:scaled-test}
 M_s(X,Z,y)\succeq0
 \quad\Longleftrightarrow\quad Z\succeq\Phi_y(X^{-1}),
 \qquad s\ne0.
\end{equation}
As $s\to0$, these matrices converge to
\[
 M(X,Z,y)=
 \begin{pmatrix}D(X)&B(y)\\B(y)^\top&E(Z)\end{pmatrix}.
\]
It follows immediately that $Z\succeq\Phi_y(X^{-1})$ implies
$M(X,Z,y)\succeq0$.

\emph{Recovering the inequality from the limit.}
For the converse, suppose that $M(X,Z,y)\succeq0$.  Positivity and
unitality give
\[
 D(X)\succeq\lambda_{\min}(X)I_a\succ0.
\]
For every $\delta>0$,
\[
 M(X,Z+\delta I_4,y)
 =M(X,Z,y)+\operatorname{diag}(0,\delta I_c)
 \succ0.
\]
To see strict positivity, a vector with nonzero lower component has
strictly positive contribution from the added term.  A nonzero
vector with lower component zero has strictly positive quadratic
form under $D(X)$.

Consequently, for each fixed $\delta>0$, the matrices
$M_s(X,Z+\delta I_4,y)$ are positive definite for sufficiently small
nonzero $s$.  Equation~\eqref{eq:scaled-test} yields
$Z+\delta I_4\succeq\Phi_y(X^{-1})$.  Letting $\delta\downarrow0$
proves the converse.  We have therefore shown
\[
 Z\succeq\Phi_y(X^{-1})
 \quad\Longleftrightarrow\quad M(X,Z,y)\succeq0.
\]
Taking the Schur complement of the positive definite block $D(X)$
gives~\eqref{eq:schur-equivalence}.
\end{proof}

The maps and their finite dimensions are now fixed by the
hypothetical pencil.  We next recover the scalar form $b$ from
\eqref{eq:schur-equivalence}; this will force a nonzero bilinear square
dominated by $b$ and contradict its defining obstruction.
